Dreiecksungleichung \qquad \qquad
\begin{math}\begin{array}{l}
		\abs{x + y} \le \abs{x} + \abs{y} \\
		\abs{\abs{x}- \abs{y}} \le \abs{x-y}
	\end{array}\end{math} \\
Arithmetische Summenformel
\begin{math}\begin{array}{l}
		\sum\limits_{k = 1}^{n}k = \frac{n (n+1)}{2}
	\end{array}\end{math}  \\
\\
Geometrische Summenformel 
\begin{math}\begin{array}{l}
		\sum\limits_{k = 0}^{n}q^k = \frac{1 - q^{n+1}}{1-q}
	\end{array}\end{math}\\
\\
Bernoulli-Ungleichung \qquad \qquad
\begin{math}\begin{array}{l}
		(1+a)^n \ge 1 + na
	\end{array}\end{math}\\
\\
Binomialkoeffizient \qquad \qquad
\begin{math}\begin{array}{l}
		\binom{n}{k} = \frac{n!}{k!(n-k)!} \\
		\binom{n}{0} = \binom{n}{n} = 1
	\end{array}\end{math}\\
\\
Binomialidentität \qquad \qquad \quad
\begin{math}\begin{array}{l}
		\binom{n}{k-1}+\binom{n}{k} = \binom{n+1}{k}
	\end{array}\end{math}\\
\\
Binomische Formel \qquad \qquad
\begin{math}\begin{array}{l}
		(a+b)^n = \sum\limits_{k = 0}^{n} \binom{n}{k} a^{n-k} b^{k}
	\end{array}\end{math}

\subsubsection{Komplexe Zahlen - Polarkoordinaten}
$z=a+b\mathbf{i}\ne0$\ in Polarkoordinaten:\\
$z=r (\cos(\varphi)+\mathbf{i}\sin(\varphi))=r\cdot e^{\mathbf{i} \varphi}$\\
$r=|z|=\sqrt{a^2+b^2}\quad\varphi=\arg(z)=\begin{cases}+\arccos \left( \frac{a}{r}\right),  & b\ge0   \\  -\arccos \left( \frac{a}{r}\right), & b<0\end{cases}$

TODO: Mengensachen, offen, geschlossen, kompakt


\subsubsection{Mengen}
\begin{description}
	\item[kompakt] Menge $\Omega$, falls jede Folge von Punkten eine konvergente Teilfolge bezitzt. Genau dann, wenn die Menge beschränkt und abgeschlossen ist
\end{description}